% header-includes.tex — Raw LaTeX injected into every PDF build.
%
% This file augments code-block padding for readability.
% No special-character workarounds — content uses ASCII (-> instead of ⇒)
% so all characters are guaranteed to render in every PDF viewer.

\usepackage{xcolor}
% Load fvextra explicitly — provides the `breaklines` option for code blocks.
% Eisvogel normally loads this conditionally (only when it sees code blocks),
% but the conditional doesn't fire reliably with stdin input, so we load it
% unconditionally here.
\usepackage{fvextra}

% ============================================================================
% CODE BLOCK STYLING — extra padding for readability
% ============================================================================
% The Highlighting environment is defined by Eisvogel using fvextra's
% Verbatim. We add extra left padding (so commands have breathing room)
% and extra bottom padding (so the code doesn't crowd surrounding text).
%
% breaklines=true wraps long lines at the margin with a small indent on
% continuation lines so wrapped URLs and commands remain readable.

\DefineVerbatimEnvironment{Highlighting}{Verbatim}{%
  breaklines,%
  fontsize=\small,%
  commandchars=\\\{\},%
  xleftmargin=1.5em,%
  framesep=0.7em,%
  baselinestretch=1.05,%
  rulecolor=\color{table-rule-color}%
}

% Extra vertical space BEFORE and AFTER each code block.
\AtBeginEnvironment{Highlighting}{\par\vspace{0.4em}}
\AtEndEnvironment{Highlighting}{\par\vspace{0.8em}}
